\section{A short return score cancels the unknown event rarity}
\label{sec:return}
The remainder uses the singleton $b$. Replacing it with an arbitrary event set would generally require a conditional distribution and a matrix-valued return interface; no such extension is assumed.

\subsection{The exact canceled quantity}
For $P=P_h^u$, $p=\pi(b)$, define
\begin{equation}
 Z(x)=\sum_{j\ge0}(P^j(x,b)-p),\qquad S_u=Z(b).
 \label{eq:green}
\end{equation}
The Poisson equation $(I-P)Z=\one_{\{b\}}-p$ and $\sum_x\pi(x)Z(x)=0$ imply
\begin{equation}
 \E_x T_b=\frac{S_u-Z(x)}p,\qquad
 M_{h,u}=\frac hN\frac{S_u}{p(1-p)}.
 \label{eq:cancel}
\end{equation}
The enormous unknown factor $1/[p(1-p)]$ is exactly common to all items. Ranking $S_u$ ranks their substitute means. This identity is established Green-function theory; the finite experimental reduction is the operation constructed here.

Since full conditionals are at most $\bar q$, the same is true under partial conditioning. Multiplying a chain of conditional probabilities, and separately multiplying odds along bit flips, gives
\begin{equation}
 p\le\bar p:=\bar q^N,\qquad
 \frac{\pi(b)}{\pi(x)}\le a^{w(x)},\quad
 a=\frac{\bar q}{1-\bar q}<1,\quad w(x)=\#\{i:x_i=0\}.
 \label{eq:massratio}
\end{equation}
No partition function or numerical weight table is required. Put
\begin{equation}
 d_-=1-e^{-r_-h(1-\omega)},\quad d_+=1-e^{-r_+h(1+\omega)},\quad
 s_0=\frac{(1-\bar p)^2}{d_+}.
 \label{eq:s0}
\end{equation}
The PSD spectral representation and Cauchy--Schwarz give
\begin{equation}
 S_u\ge\frac{(1-p)^2}{1-P(b,b)}\ge s_0.
 \label{eq:greenlower}
\end{equation}
The nonconstant spectral weights sum to $1-p$, their weighted $1-\lambda$ sum is $1-P(b,b)$, and the latter is at most $d_+$.

\subsection{Stop after a few recoveries}
For $1\le k<N$ let $C_k=\{x:w(x)\ge k\}$ and let $\sigma_k$ be its first entrance time from $b$. Count
\begin{equation}
 V_u=\sum_{j=0}^{\sigma_k-1}\one_{\{X_j=b\}},\qquad A_u=\E_b V_u.
 \label{eq:score}
\end{equation}
Each visit to $b$, including a self-loop, restarts the same chance of reaching $C_k$ before the next visit. Therefore $V_u$ is geometric on $\{1,2,\ldots\}$ and
\begin{equation}
 A_u=\frac1{\Pr_b(\sigma_k<T_b^+)},\qquad
 \operatorname{Var}(V_u)=A_u(A_u-1)\le A_u^2.
 \label{eq:geometric}
\end{equation}
Equivalently $A_u=\pi(b)/\operatorname{cap}_P(b,C_k)$. Capacity and renewal theory supply the representation; the next bounds make this unknown-system measurement finite and short.

While $w<k$ its drift is at least
\begin{equation}
 \nu=\left(1-\frac kN\right)d_-(1-\bar q)-\frac kN d_+\bar q.
 \label{eq:drift}
\end{equation}
Require $\nu>0$. Optional stopping and the bounded-increment martingale inequality yield
\begin{equation}
 \E_b\sigma_k\le k/\nu,\qquad
 \Pr_b(\sigma_k>J)\le e^{-(\nu J-k)^2/(2J)},\quad J\ge k/\nu.
 \label{eq:stopcost}
\end{equation}
For the second assertion extend the stopped count by deterministic drift $\nu$ after exit. Its increments have range length two and conditional drift at least $\nu$. On $\sigma_k>J$ it is still below $k$. If $J\ge2k/\nu$ the displayed bound is at most $e^{-\nu^2J/8}$.

\begin{theorem}[Finite local-return certificate]\label{thm:localreturn}
For any positive integer analysis horizon $L$, set
\begin{equation}
 D=\bar p(L+k/\nu)+La^k+\frac{a^{k/2}e^{-gL}}g.
 \label{eq:D}
\end{equation}
Then, uniformly over the menu, $|S_u-A_u|\le D$. The horizon $L$ is used only in certificate arithmetic, not as the running length of the probe.
\end{theorem}
\begin{proof}
For a boundary state $x$ with $w(x)=k$, reversibility gives $P^j(x,b)\le\pi(b)/\pi(x)\le a^k$. The first $L$ terms of $Z(x)$ are between $-Lp$ and $La^k$. Spectral contraction gives
\[
 |P^j(x,b)-p|\le\sqrt{\pi(b)/\pi(x)}(1-g)^j\le a^{k/2}(1-g)^j,
\]
so the remaining sum has modulus at most $a^{k/2}e^{-gL}/g$. The stopped Poisson identity is
$S_u=A_u-p\E_b\sigma_k+\E_bZ(X_{\sigma_k})$.
Apply these bounds and~\eqref{eq:stopcost}. The finite state space and positive drift justify stopping.
\end{proof}

For $0<\eta_b<1/4$, one finite schedule is
\begin{align}
 L&=\max\left\{1,\left\lceil g^{-1}\log\max\left(1,\frac3{\eta_bgs_0}\right)\right\rceil\right\},\\
 k&=\max\left\{1,\left\lceil\frac{\log\max(1,3L/(\eta_bs_0))}{\log(1/a)}\right\rceil\right\}.
 \label{eq:schedule}
\end{align}
Check $k<N$, $\nu>0$, $\bar p(L+k/\nu)\le\eta_b s_0/3$, and directly check $D\le\eta_b s_0$ with outward rounding. At fixed margins these checks succeed for all sufficiently large $N$: $g^{-1}=O(N)$, $L=O(N\log N)$, $k=O(\log N)$, and $\bar p$ decreases exponentially. The threshold is computed. A boundary without certified slack triggers refinement or the fallback below.

\subsection{Confidence, finite caps and ranking}
A geometric variable of mean $A\ge1$ obeys
$\E e^{tV}\le(1-At)^{-1}$ for $0\le t<1/A$ and $\E V^2\le2A^2$.
For the first bound use its exact generating function and $(1-t)e^t\le1$; for the lower tail use $e^{-x}\le1-x+x^2/2$. Chernoff's argument gives
\begin{equation}
 \Pr\big(|\widehat A_u/A_u-1|>\eta_s\big)\le2e^{-n\eta_s^2/4},
 \quad0<\eta_s\le1.
 \label{eq:concentration}
\end{equation}
The conditional version follows by iteration of conditional moment bounds; no independence of reused physical cells is presumed beyond the declared conditional interface.
Use
\begin{equation}
 n=\left\lceil4\eta_s^{-2}\log\frac{2K}{\delta_{\rm stat}}\right\rceil,
 \quad J=\left\lceil\max\left\{2k/\nu,\,8\nu^{-2}\log\frac{Kn}{\delta_{\rm cap}}\right\}\right\rceil.
 \label{eq:finitecaps}
\end{equation}
Abort explicitly if any probe fails to reach $C_k$ by $J$ updates. Its ideal probability is at most $\delta_{\rm cap}$. Otherwise the finite algorithm agrees with the uncapped reference. Choosing the largest sample mean then gives, with confidence $1-\delta_{\rm stat}-\delta_{\rm cap}$,
\begin{equation}
 \frac{M_{\widehat u}}{\max_uM_u}
 \ge\frac{\kappa_-}{\kappa_+}
 \frac{(1-\eta_s)(1-\eta_b)}{(1+\eta_s)(1+\eta_b)}.
 \label{eq:decisionfactor}
\end{equation}
Here $D\le\eta_b s_0\le\eta_b S_u$ and the common $p$ cancels.

The observed certificate can be sharper. Set
\begin{equation}
 S_u^- =\max\{s_0,\widehat A_u/(1+\eta_s)-D\},\quad
 S_u^+ =\widehat A_u/(1-\eta_s)+D.
 \label{eq:observedinterval}
\end{equation}
Then $M_u/M_v\ge(\kappa_-/\kappa_+)S_u^-/S_v^+$. If the latter exceeds one for every competitor, the protocol certifies the actual best mean-lifetime item while leaving the absolute lifetime unknown. Adaptive sample stages must divide their confidence budget; repeated use of a fixed-time interval is invalid.

\subsection{A finite fallback without rarity reconstruction}
If the sparse certificate fails, use full length-$L$ counts $V_L=\sum_{j<L}\one_{\{X_j=b\}}$. Their means are $S_u+Lp-T_u$, where $0\le T_u\le e^{-gL}/g$. Choose $L$ so that this last bound is at most $\eta s_0$, and estimate each count mean to additive $\eta s_0$ using
\[
 n_{\rm fb}=\left\lceil\frac{L^2}{2\eta^2s_0^2}\log\frac{2K}{\delta_{\rm stat}}\right\rceil.
\]
The common unknown $Lp$ cancels in comparisons. The selected $S$ is at least the best minus $3\eta s_0$, hence within relative factor $1-3\eta$. This is polynomial in $N$, numerical inverse margins and accuracy, not in $1/p$. At small $N$ an ordinary exact solver can be better; this fallback is not claimed optimal.
